\subsection{Definition of associated prime ideals}

DEFINITION 1. Let M be a module over a ring A. A prime ideal p is said to be associated with M if there exists \(x \in M\) such that p is equal to the annihilator of x. The set of prime ideals associated with M is denoted by \(\operatorname{Ass}_{\mathbf{A}}(\mathbf{M})\) , or simply \(\operatorname{Ass}(\mathbf{M})\) .

* Example. Let \(a\) be an ideal in the polynomial ring \(A = \mathbf{C}[X_1, \ldots, X_r]\) , \(V\) the corresponding affine algebraic variety and \(V_1, \ldots, V_p\) the irreducible components of \(V\) . If \(M\) is taken to be the ring \(A/a\) of functions which are regular on \(V\) , the set of prime ideals associated with \(M\) consists of the ideals of \(V_1, \ldots, V_p\) and in general other prime ideals each of which contains one of the ideals of the \(V_1\) .

As the annihilator of 0 is A, an element \(x \in \mathbf{M}\) whose annihilator is a prime ideal is necessarily \(\neq 0\) . To say that a prime ideal \(p\) is associated with M amounts to saying that M contains a submodule isomorphic to A/p (namely Ax, for all \(x \in M\) whose annihilator is p).

If an A-module M is the union of a family \((M_{t})_{t\in I}\) of submodules, then clearly

\[
\operatorname{Ass} (\mathbf {M}) = \bigcup_ {\iota \in I} \operatorname{Ass} \left(\mathbf {M} _ {\iota}\right).\tag{1}
\]

PROPOSITION 1. For every prime ideal \(\mathfrak{p}\) of a ring \(\mathbf{A}\) and every submodule \(\mathbf{M} \neq 0\) of \(\mathbf{A} / \mathfrak{p}\) , \(\operatorname{Ass}(\mathbf{M}) = \{\mathfrak{p}\}\) .

As the ring A/p is an integral domain, the annihilator of an element \(\neq0\) of A/p is p.

PROPOSITION 2. Let M be a module over a ring A. Every maximal element of the set of ideals Ann(x) of A, where x runs through the set of elements \(\neq0\) of M, belongs to Ass(M).

Let \(\mathfrak{a} = \operatorname{Ann}(x)\) ( \(x \in \mathbf{M}, x \neq 0\) ) be such a maximal element; it is sufficient to show that \(\mathfrak{a}\) is prime. As \(x \neq 0\) , \(\mathfrak{a} \neq \mathbf{A}\) . Let \(b, c\) be elements of \(\mathbf{A}\) such that \(bc \in \mathfrak{a}\) and \(c \notin \mathfrak{a}\) . Then \(cx \neq 0\) , \(b \in \operatorname{Ann}(cx)\) and \(\mathfrak{a} \subset \operatorname{Ann}(cx)\) . As \(\mathfrak{a}\) is maximal, \(\operatorname{Ann}(cx) = \mathfrak{a}\) , whence \(b \in \mathfrak{a}\) , so that \(\mathfrak{a}\) is prime.

COROLLARY 1. Let M be a module over a Noetherian ring A. Then the condition \(M \neq \{0\}\) is equivalent to \(\operatorname{Ass}(M) \neq \varnothing\) .

If \(M = \{0\}\) , clearly \(\text{Ass}(M)\) is empty (without any hypothesis on A). If \(M \neq \{0\}\) , the set of ideals of the form \(\text{Ann}(x)\) , where \(x \in M\) and \(x \neq 0\) , is non-empty and consists of ideals \(\neq A\) ; as A is Noetherian, this set has a maximal element; then it suffices to apply Proposition 2.

COROLLARY 2. Let A be a Noetherian ring, M an A-module and a an element of A. For the homothety on M with ratio a to be injective, it is necessary and sufficient that a belong to no prime ideal associated with M.

If \(a\) belongs to a prime ideal \(\mathfrak{p} \in \operatorname{Ass}(\mathbf{M})\) , then \(\mathfrak{p} = \operatorname{Ann}(x)\) , where \(x \in \mathbf{M}\) , \(x \neq 0\) ; whence \(ax = 0\) and the homothety with ratio \(a\) is not injective. Conversely, if \(ax = 0\) for some \(x \in \mathbf{M}\) such that \(x \neq 0\) , then \(Ax \neq \{0\}\) , whence \(\operatorname{Ass}(x) \neq \varnothing\) (Corollary 1). Let \(\mathfrak{p} \in \operatorname{Ass}(\mathbf{A}x)\) ; then obviously \(\mathfrak{p} \in \operatorname{Ass}(\mathbf{M})\) and \(\mathfrak{p} = \operatorname{Ann}(bx)\) , where \(b \in \mathbf{A}\) ; whence \(a \in \mathfrak{p}\) , since \(abx = 0\) .

COROLLARY 3. The set of divisors of zero in a Noetherian ring A is the union of the ideals \(\mathfrak{p} \in \operatorname{Ass}(A)\) .

PROPOSITION 3. Let A be a ring, M an A-module and N a submodule of N. Then

\[
\operatorname{Ass} (\mathbf {N}) \subset \operatorname{Ass} (\mathbf {M}) \subset \operatorname{Ass} (\mathbf {N}) \cup \operatorname{Ass} (\mathbf {M} / \mathbf {N}).\tag{2}
\]

The inclusion \(\operatorname{Ass}(\mathbf{N}) \subset \operatorname{Ass}(\mathbf{M})\) is obvious. Let \(\mathfrak{p} \in \operatorname{Ass}(\mathbf{M})\) , \(\mathbf{E}\) be a submodule of \(\mathbf{M}\) isomorphic to \(\mathbf{A} / \mathfrak{p}\) and \(\mathbf{F} = \mathbf{E} \cap \mathbf{N}\) . If \(\mathbf{F} = \{0\}\) , \(\mathbf{E}\) is isomorphic to a submodule of M/N, whence \(\mathfrak{p} \in \operatorname{Ass}(M/N)\) . If \(F \neq \{0\}\) , the annihilator of every element \(\neq 0\) of F is p (Proposition 1) and hence \(\mathfrak{p} \in \operatorname{Ass}(F) \subset \operatorname{Ass}(N)\) .

COROLLARY 1. If an A-module M is the direct sum of a family \((\mathbf{M}_t)_{t \in I}\) of submodules, then \(\operatorname{Ass}(\mathbf{M}) = \bigcup_{t \in I} \operatorname{Ass}(\mathbf{M}_t)\) .

It may be reduced to the case where I is finite by means of (1), then to the case where \(\operatorname{Card}(I) = 2\) by induction on \(\operatorname{Card}(I)\) . Then let \(I = \{i, j\}\) , where \(i \neq j\) ; as \(M/M_{i}\) is isomorphic to \(M_{j}\) , \(\operatorname{Ass}(M) \subset \operatorname{Ass}(M_{i}) \cup \operatorname{Ass}(M_{j})\) (Proposition 3); moreover, \(\operatorname{Ass}(M_{i})\) and \(\operatorname{Ass}(M_{j})\) are contained in \(\operatorname{Ass}(M)\) (Proposition 3), whence the result.

COROLLARY 2. Let M be an A-module and \((\mathbf{Q}_i)_{i\in I}\) a finite family of submodules of M. If \(\bigcap_{i\in I} \mathbf{Q}_i = \{0\}\) , then

\[
\operatorname{Ass} (\mathbf {M}) \subset \bigcup_ {i \in I} \operatorname{Ass} (\mathbf {M} / \mathbf {Q} _ {i}).
\]

The canonical mapping \(\mathbf{M} \to \bigoplus_{i \in I} (\mathbf{M} / \mathbf{Q}_i)\) is injective; then it suffices to apply Proposition 3 and its Corollary 1.

PROPOSITION 4. Let M be an A-module and \(\Phi\) a subset of \(\operatorname{Ass}(\mathbf{M})\) . Then there exists a submodule N of M such that \(\operatorname{Ass}(\mathbf{N}) = \operatorname{Ass}(\mathbf{M}) - \Phi\) and \(\operatorname{Ass}(\mathbf{M}/\mathbf{N}) = \Phi\) .

Let \(\mathfrak{E}\) be the set of submodules \(P\) of \(M\) such that \(\operatorname{Ass}(P) \subset \operatorname{Ass}(M) - \Phi\) . Formula (1) shows that the set \(\mathfrak{E}\) , ordered by inclusion, is inductive; moreover, \(\{0\} \in \mathfrak{E}\) and hence \(\mathfrak{E} \neq \varnothing\) . Let \(N\) be a maximal element of \(\mathfrak{E}\) . Then \(\operatorname{Ass}(N) \subset \operatorname{Ass}(M) - \Phi\) . We shall see that \(\operatorname{Ass}(M/N) \subset \Phi\) , which, by Proposition 3, will complete the proof. Let \(p \in \operatorname{Ass}(M/N)\) ; then \(M/N\) contains a submodule \(F/N\) isomorphic to \(A/p\) . By Propositions 1 and 3, \(\operatorname{Ass}(F) \subset \operatorname{Ass}(N) \cup \{p\}\) . Since \(N\) is maximal in \(\mathfrak{E}\) , \(F \notin \mathfrak{E}\) and hence \(p \in \Phi\) .

